\chapter{The Fundamental Group and Coverings}

\begin{center}\itshape
Walk around a circle and you come back to where you started, but you have gone around once. Walk around twice and you have gone around twice, and no
amount of wiggling turns one walk into the other. Loops up to deformation form a group, and that group sees holes. For the circle it is the integers,
which is the winding number of Volume~7 in a new coat; for the figure eight it is a free group; and the subgroups of a free group, which look
hopelessly complicated, turn out to be free again, because a covering of a graph is a graph.
\end{center}

In 1895 Henri Poincar\'e published a memoir of more than a hundred pages in the \emph{Journal de l'\'Ecole Polytechnique} under the title
\emph{Analysis situs}, the old name for what we now call topology. He wanted to tell manifolds apart, and he gave two tools for doing so. One was a
count of independent cycles, the Betti numbers, which Chapter~6 turns into homology. The other was new: the group of loops, which he called the
\emph{groupe fondamental}. Over the next nine years he published five supplements to correct and extend the memoir, and they are a rare public record of
a great mathematician getting things wrong and fixing them.

In the second supplement Poincar\'e had stated that a three-dimensional manifold with the Betti numbers of the sphere is a sphere. In the fifth, in
1904, he constructed a counterexample to his own claim: a manifold, now called the Poincar\'e homology sphere, which has the same homology as the
three-dimensional sphere but a nontrivial fundamental group (it has $120$ elements). He then asked whether a closed three-manifold whose fundamental group is
trivial must be a sphere, and stopped there, remarking that the question would take him too far. It took the rest of the century.

The answer came in November 2002, when Grigori Perelman, a researcher at the Steklov Institute in Saint Petersburg, posted the first of three preprints on
the arXiv. They completed a programme of Richard Hamilton, which deforms the metric of a manifold by its curvature until its shape becomes visible. Several
teams spent the next years checking the details. In August 2006 the International Congress of Mathematicians in Madrid awarded Perelman a Fields Medal,
which he declined, the first person to do so; in 2010 he also declined the million-dollar prize the Clay Mathematics Institute had offered for the
problem. Poincar\'e's question is the only one of the Clay Millennium Problems solved to date.

This chapter builds Poincar\'e's group, computes it for circles, spheres, tori, graphs and surfaces, and develops its partner, the theory of covering
spaces, in which subgroups of the fundamental group correspond to coverings exactly as subgroups of a Galois group correspond to fields. Throughout,
spaces are metric spaces (all our examples are), maps are continuous, and $I=[0,1]$.

\section{Loops and the fundamental group}

\begin{definition}
A \emph{path} in $X$ is a map $\gamma\colon I\to X$; it goes from $\gamma(0)$ to $\gamma(1)$, and it is a \emph{loop} at $x_0$ if $\gamma(0)=\gamma(1)=x_0$. Two paths $\gamma_0,\gamma_1$ with the same
end points are \emph{homotopic}, written $\gamma_0\simeq\gamma_1$, if there is a map $H\colon I\times I\to X$ with $H(t,0)=\gamma_0(t)$, $H(t,1)=\gamma_1(t)$, and $H(0,s)$, $H(1,s)$ independent of $s$
(the end points stay fixed). If $\gamma(1)=\delta(0)$, the \emph{product} $\gamma\cdot\delta$ runs through $\gamma$ at double speed, then through $\delta$: it is $\gamma(2t)$ for $t\le\frac12$ and
$\delta(2t-1)$ for $t\ge\frac12$. The \emph{reverse} of $\gamma$ is $\bar\gamma(t)=\gamma(1-t)$, and $c_x$ denotes the constant path at $x$.
\end{definition}

Homotopy is an equivalence relation: reverse a homotopy for symmetry, and for transitivity run two homotopies one after the other, each at double
speed in $s$ (the result is continuous because the two pieces agree on the line $s=\frac12$, a closed set). We write $[\gamma]$ for the homotopy class.

\begin{lemma}\label{lem:convexhomotopy}
\begin{enumerate}
\item[(a)] Two paths with the same end points in a convex subset of $\R^n$ are homotopic, by the straight-line homotopy $H(t,s)=(1-s)\gamma_0(t)+s\gamma_1(t)$.
\item[(b)] If $f\colon X\to Y$ is a map and $\gamma_0\simeq\gamma_1$ in $X$, then $f\circ\gamma_0\simeq f\circ\gamma_1$ in $Y$.
\item[(c)] (Reparametrization.) If $\varphi\colon I\to I$ is continuous with $\varphi(0)=0$ and $\varphi(1)=1$, then $\gamma\circ\varphi\simeq\gamma$.
\end{enumerate}
\end{lemma}

\begin{proof}
(a) The segment between two points of a convex set stays in it, and $H(0,s)$, $H(1,s)$ are constant because $\gamma_0$ and $\gamma_1$ have the same end points. (b) Compose the
homotopy with $f$. (c) The paths $\varphi$ and $\mathrm{id}_I$ are homotopic in the convex set $I$ by (a); apply (b) with $f=\gamma$.
\end{proof}

\begin{proposition}\label{prop:pi1group}
Whenever the products are defined:
\begin{enumerate}
\item[(a)] if $\gamma\simeq\gamma'$ and $\delta\simeq\delta'$ then $\gamma\cdot\delta\simeq\gamma'\cdot\delta'$;
\item[(b)] $(\gamma\cdot\delta)\cdot\varepsilon\simeq\gamma\cdot(\delta\cdot\varepsilon)$;
\item[(c)] $c_{\gamma(0)}\cdot\gamma\simeq\gamma\simeq\gamma\cdot c_{\gamma(1)}$;
\item[(d)] $\gamma\cdot\bar\gamma\simeq c_{\gamma(0)}$.
\end{enumerate}
Consequently the homotopy classes of loops at $x_0$ form a group under $[\gamma][\delta]=[\gamma\cdot\delta]$, with identity $[c_{x_0}]$ and inverse $[\gamma]^{-1}=[\bar\gamma]$: the
\emph{fundamental group} $\pi_1(X,x_0)$.
\end{proposition}

\begin{proof}
(a) Put the two homotopies side by side, each at double speed in $t$; they agree at $t=\frac12$ because the end points are fixed. (b) The left side is the right side
reparametrized by the piecewise linear $\varphi$ sending $0,\frac14,\frac12,1$ to $0,\frac12,\frac34,1$; apply Lemma~\ref{lem:convexhomotopy}(c). (c) Likewise $c\cdot\gamma=\gamma\circ\varphi$ with
$\varphi=0$ on $[0,\frac12]$ and $\varphi(t)=2t-1$ after, and similarly on the right. (d) $H(t,s)=\gamma(2t(1-s))$ for $t\le\frac12$ and $H(t,s)=\gamma((2-2t)(1-s))$ for $t\ge\frac12$
runs out along $\gamma$ only as far as $\gamma(1-s)$ and comes back; it is $\gamma\cdot\bar\gamma$ at $s=0$ and constant at $s=1$.
\end{proof}

A path $h$ from $x_0$ to $x_1$ gives a map $\beta_h\colon\pi_1(X,x_0)\to\pi_1(X,x_1)$, $[\gamma]\mapsto[\bar h\cdot\gamma\cdot h]$: go back along $h$, do the loop, return. By the proposition it is a
homomorphism ($\bar h\gamma h\bar h\delta h\simeq\bar h\gamma\delta h$) with inverse $\beta_{\bar h}$. So in a path-connected space the fundamental group does not depend on the base
point, up to isomorphism, and we often write $\pi_1(X)$. The space is \emph{simply connected} if it is path-connected and $\pi_1(X)$ is trivial. A convex subset of $\R^n$ is
simply connected by Lemma~\ref{lem:convexhomotopy}(a).

A map $f\colon X\to Y$ induces $f_*\colon\pi_1(X,x_0)\to\pi_1(Y,f(x_0))$, $[\gamma]\mapsto[f\circ\gamma]$, well defined by Lemma~\ref{lem:convexhomotopy}(b) and a homomorphism because
$f\circ(\gamma\cdot\delta)=(f\circ\gamma)\cdot(f\circ\delta)$. Clearly $(g\circ f)_*=g_*\circ f_*$ and $\mathrm{id}_*=\mathrm{id}$, so a homeomorphism induces an isomorphism.

\begin{proposition}\label{prop:homotopyinv}
Let $F\colon X\times I\to Y$ be a homotopy from $f=F(\cdot,0)$ to $g=F(\cdot,1)$, and let $h(s)=F(x_0,s)$. Then $g_*=\beta_h\circ f_*$. Consequently, if $f\colon X\to Y$ is a
\emph{homotopy equivalence}, that is, if there is $g\colon Y\to X$ with $g\circ f$ and $f\circ g$ homotopic to the identities, then $f_*$ is an isomorphism.
\end{proposition}

\begin{proof}
For a loop $\gamma$ at $x_0$, consider $G(t,s)=F(\gamma(t),s)$ on the square. In the square, the bottom edge is homotopic, with end points fixed, to the path that goes up the left side,
along the top and down the right side (Lemma~\ref{lem:convexhomotopy}(a)). Composing with $G$ (part (b)), $f\circ\gamma\simeq h\cdot(g\circ\gamma)\cdot\bar h$, that is,
$[g\circ\gamma]=[\bar h\cdot(f\circ\gamma)\cdot h]$.

For the second statement, $(g\circ f)_*=g_*\circ f_*$ is the composite of the identity with some $\beta$, an isomorphism; so $f_*$ is injective and $g_*$ (on $\pi_1(Y,f(x_0))$) is
surjective. Applying the same argument to $f\circ g$ at the base point $f(x_0)$ shows that $g_*$ is also injective. So $g_*$ is bijective, and then so is $f_*$.
\end{proof}

A subspace $A\subset X$ is a \emph{deformation retract} of $X$ if there is a homotopy $F\colon X\times I\to X$ from the identity to a map into $A$, with $F(a,s)=a$ for $a\in A$ and all $s$.
Then the inclusion $A\to X$ is a homotopy equivalence (its homotopy inverse is $F(\cdot,1)$), so $\pi_1(A,a)\cong\pi_1(X,a)$. For example, $\R^2\setminus\{0\}$ deformation retracts onto the
unit circle by $F(x,s)=(1-s)x+s\frac x{|x|}$, and a solid torus onto its core circle.

\section{Covering spaces and the circle}

\begin{definition}
A \emph{covering} is a map $p\colon\tilde X\to X$ such that every point of $X$ has an open neighbourhood $U$ that is \emph{evenly covered}: $p^{-1}(U)$ is a disjoint union of open
sets, the \emph{sheets} over $U$, each mapped homeomorphically onto $U$ by $p$. A \emph{lift} of a map $f\colon Y\to X$ is a map $\tilde f\colon Y\to\tilde X$ with $p\circ\tilde f=f$.
\end{definition}

\begin{example}\label{ex:coverings}
(a) $p\colon\R\to S^1$, $p(t)=e^{2\pi it}$: an open arc $U$ of $S^1$ other than the whole circle is evenly covered, its preimage being a disjoint union of translates of an interval. Above
each point sits a copy of $\Z$; the picture is a helix over the circle (Figure~\ref{fig:helix}). (b) $z\mapsto z^n$ from $S^1$ to itself, with $n$ sheets. (c) $\R^2\to T^2=\R^2/\Z^2$.
(d) The quotient map $S^n\to\R P^n$ that identifies antipodal points: an open set of $S^n$ contained in an open hemisphere maps homeomorphically onto its image, and
the preimage of that image is the set and its antipode, two disjoint sheets.
\end{example}

\begin{figure}[ht]
\centering
\fig{ch05-helix}
\caption{The covering $\R\to S^1$ drawn as a helix over the circle. A loop that goes twice around the circle lifts to a path that climbs two turns of the helix: the
height it gains is the number of turns, and no homotopy of the loop can change it.}
\label{fig:helix}
\end{figure}

\begin{lemma}[uniqueness of lifts]\label{lem:uniquelift}
Let $p\colon\tilde X\to X$ be a covering and $Y$ connected. Two lifts of the same map $f\colon Y\to X$ that agree at one point agree everywhere.
\end{lemma}

\begin{proof}
Let $\tilde f_1,\tilde f_2$ be the lifts. For $y\in Y$, let $U$ be an evenly covered neighbourhood of $f(y)$, and $\tilde U_1,\tilde U_2$ the sheets containing $\tilde f_1(y)$, $\tilde f_2(y)$. On the
neighbourhood $W=\tilde f_1^{-1}(\tilde U_1)\cap\tilde f_2^{-1}(\tilde U_2)$ of $y$: if $\tilde U_1=\tilde U_2$, then $\tilde f_1=\tilde f_2$ on $W$, because $p$ is injective on a sheet; if $\tilde U_1\ne\tilde U_2$, the sheets are
disjoint and $\tilde f_1\ne\tilde f_2$ on $W$. So the set where the lifts agree and the set where they differ are both open. The first is nonempty, and $Y$ is connected.
\end{proof}

\begin{theorem}[homotopy lifting]\label{thm:homotopylifting}
Let $p\colon\tilde X\to X$ be a covering, $F\colon I\times I\to X$ a map, and $\tilde x$ a point with $p(\tilde x)=F(0,0)$. There is a unique lift $\tilde F$ of $F$ with $\tilde F(0,0)=\tilde x$. In particular
(taking $F(t,s)=\gamma(t)$) every path $\gamma$ has a unique lift $\tilde\gamma$ starting at any given point over $\gamma(0)$.
\end{theorem}

\begin{proof}
Uniqueness is Lemma~\ref{lem:uniquelift}. The open sets $F^{-1}(U)$, $U$ evenly covered, cover the compact square, so they have a Lebesgue number (Volume~6, Theorem~4.8): we can cut
the square into small squares $R_{ij}=[\frac{i-1}N,\frac iN]\times[\frac{j-1}N,\frac jN]$, each mapped by $F$ into an evenly covered set $U_{ij}$. Lift them one at a time, in the order
$R_{11},R_{21},\dots,R_{N1},R_{12},\dots$, row by row. When we reach $R_{ij}$, the lift is already defined on the set $L$ made of its bottom edge and its left edge (on the parts that
belong to earlier squares; for $R_{11}$, on the single point $(0,0)$). The set $L$ is connected and $\tilde F(L)$ lies in $p^{-1}(U_{ij})$, so it lies in a single sheet $\tilde U$ (the sheets are
open, disjoint, and $\tilde F(L)$ is connected). Define $\tilde F=(p|_{\tilde U})^{-1}\circ F$ on $R_{ij}$. This agrees with the earlier definition on $L$, since both are points of $\tilde U$ over the same
point of $X$; so the pieces glue into a map continuous on each closed square, hence continuous.
\end{proof}

\begin{corollary}\label{cor:liftendpoints}
If $\gamma_0\simeq\gamma_1$ and $\tilde\gamma_0,\tilde\gamma_1$ are their lifts from the same point, then $\tilde\gamma_0(1)=\tilde\gamma_1(1)$ and $\tilde\gamma_0\simeq\tilde\gamma_1$. Consequently
$p_*\colon\pi_1(\tilde X,\tilde x_0)\to\pi_1(X,p(\tilde x_0))$ is injective.
\end{corollary}

\begin{proof}
Lift the homotopy $H$ from the starting point. On the left side of the square $H$ is constant, so its lift is a lift of a constant path, which is constant by uniqueness; the same holds on the
right side. On the bottom and top edges the lift is a lift of $\gamma_0$, $\gamma_1$ from $\tilde x$, hence equals $\tilde\gamma_0$, $\tilde\gamma_1$. So $\tilde H$ is a homotopy with fixed end points. If $\tilde\lambda$ is
a loop with $p\circ\tilde\lambda\simeq c$, lift that homotopy: $\tilde\lambda$ is homotopic to the lift of a constant, which is constant.
\end{proof}

\begin{theorem}\label{thm:pi1circle}
$\pi_1(S^1,1)\cong\Z$. The isomorphism sends the class of a loop $\gamma$ to $\tilde\gamma(1)$, where $\tilde\gamma$ is the lift of $\gamma$ to $\R$ starting at $0$ for the covering $p(t)=e^{2\pi it}$, and the
class of $\omega_n(t)=e^{2\pi int}$ to $n$.
\end{theorem}

\begin{proof}
Since $p(\tilde\gamma(1))=1$, $\tilde\gamma(1)\in\Z$; call it $\Phi([\gamma])$, well defined by Corollary~\ref{cor:liftendpoints}. If $m=\Phi([\gamma])$ and $\tilde\delta$ is the lift of $\delta$ from $0$, then
$\tilde\gamma\cdot(m+\tilde\delta)$ is a path from $0$ lifting $\gamma\cdot\delta$ (as $p$ is $1$-periodic), so $\Phi([\gamma][\delta])=m+\tilde\delta(1)$: $\Phi$ is a homomorphism. The lift of $\omega_n$ is $t\mapsto nt$,
so $\Phi([\omega_n])=n$ and $\Phi$ is onto. If $\Phi([\gamma])=0$, $\tilde\gamma$ is a loop in the convex set $\R$, homotopic to the constant; pushing the homotopy down by $p$ gives
$\gamma\simeq c_1$.
\end{proof}

The lift $\tilde\gamma$ is $\frac1{2\pi}$ times a continuous argument of $\gamma$ in the sense of Volume~7, Lemma~9.1, so $\Phi$ is the winding number of $\gamma$ around $0$. Theorem~\ref{thm:pi1circle}
says that the winding number is a complete invariant: two loops in $\C^*$ are homotopic exactly when they wind equally often. For a covering
$p\colon\tilde X\to X$ and $x_0\in X$, the group $\pi_1(X,x_0)$ acts on the \emph{fibre} $p^{-1}(x_0)$ on the right: for $\tilde x$ in the fibre, $\tilde x\cdot[\gamma]$ is the end point of the lift of $\gamma$
from $\tilde x$. This is well defined by Corollary~\ref{cor:liftendpoints}, and $(\tilde x\cdot[\gamma])\cdot[\delta]=\tilde x\cdot[\gamma\delta]$ because lifts concatenate.

\begin{proposition}\label{prop:fibre}
Let $p\colon\tilde X\to X$ be a covering, $\tilde X$ path-connected, $p(\tilde x_0)=x_0$ and $H=p_*\pi_1(\tilde X,\tilde x_0)$. The action of $\pi_1(X,x_0)$ on the fibre is transitive, and the stabilizer
of $\tilde x_0$ is $H$. So $\tilde x_0\cdot g\mapsto Hg$ is a bijection from the fibre to the set of right cosets of $H$, and the number of sheets is the index $[\pi_1(X,x_0):H]$.
\end{proposition}

\begin{proof}
A path $\tilde\gamma$ in $\tilde X$ from $\tilde x_0$ to $\tilde x_1$ in the fibre is the lift of the loop $p\circ\tilde\gamma$, so $\tilde x_1=\tilde x_0\cdot[p\circ\tilde\gamma]$: the action is transitive. A loop $\gamma$
has $\tilde x_0\cdot[\gamma]=\tilde x_0$ exactly when its lift is a loop $\tilde\gamma$, and then $[\gamma]=p_*[\tilde\gamma]\in H$; conversely, for $[\gamma]=p_*[\tilde\lambda]$, the lift of $p\circ\tilde\lambda$ is $\tilde\lambda$, a loop.
\end{proof}

From now on, all spaces are path-connected and \emph{locally path-connected}: every neighbourhood of a point contains a path-connected open neighbourhood. Manifolds and graphs are.

\begin{theorem}[lifting criterion]\label{thm:liftcrit}
Let $p\colon(\tilde X,\tilde x_0)\to(X,x_0)$ be a covering and $f\colon(Y,y_0)\to(X,x_0)$ a map, $Y$ path-connected and locally path-connected. A lift $\tilde f$ with $\tilde f(y_0)=\tilde x_0$ exists if and only
if $f_*\pi_1(Y,y_0)\subset p_*\pi_1(\tilde X,\tilde x_0)$, and it is then unique.
\end{theorem}

\begin{proof}
If $\tilde f$ exists, $f_*=p_*\circ\tilde f_*$ has image in $p_*\pi_1(\tilde X,\tilde x_0)$. Conversely, for $y\in Y$ choose a path $\eta$ from $y_0$ to $y$ and let $\tilde f(y)$ be the end point of the lift of
$f\circ\eta$ from $\tilde x_0$. If $\eta'$ is another such path, $[f\circ(\eta\cdot\bar\eta')]=p_*[\tilde\lambda]$ for a loop $\tilde\lambda$ at $\tilde x_0$, so by Corollary~\ref{cor:liftendpoints} the lift of
$(f\circ\eta)\cdot(f\circ\bar\eta')$ from $\tilde x_0$ ends where the lift of $p\circ\tilde\lambda$ ends, at $\tilde x_0$. That lift is the lift of $f\circ\eta$ followed by a lift of $f\circ\bar\eta'$; reading the second half backwards,
the lifts of $f\circ\eta$ and $f\circ\eta'$ from $\tilde x_0$ end at the same point. So $\tilde f$ is well defined, and $p\circ\tilde f=f$.

Continuity at $y$: take an evenly covered $U\ni f(y)$, the sheet $\tilde U$ containing $\tilde f(y)$, and a path-connected open $W\ni y$ with $f(W)\subset U$. For $y'\in W$, use the path $\eta\cdot\sigma$
with $\sigma$ a path in $W$ from $y$ to $y'$; the lift of $f\circ\sigma$ from $\tilde f(y)$ is $(p|_{\tilde U})^{-1}\circ f\circ\sigma$. Hence $\tilde f=(p|_{\tilde U})^{-1}\circ f$ on $W$, which is continuous. Uniqueness is
Lemma~\ref{lem:uniquelift}.
\end{proof}

\section{Deck transformations and the Galois correspondence}

\begin{definition}
A \emph{deck transformation} of a covering $p\colon\tilde X\to X$ is a homeomorphism $\tau$ of $\tilde X$ with $p\circ\tau=p$; they form a group $\operatorname{Deck}(p)$. Two coverings $p_1\colon\tilde X_1\to X$ and
$p_2\colon\tilde X_2\to X$ are \emph{isomorphic} if there is a homeomorphism $\phi\colon\tilde X_1\to\tilde X_2$ with $p_2\circ\phi=p_1$. A covering with $\tilde X$ path-connected is \emph{normal} if
$\operatorname{Deck}(p)$ acts transitively on each fibre, and \emph{universal} if $\tilde X$ is simply connected.
\end{definition}

A deck transformation that fixes one point is the identity (Lemma~\ref{lem:uniquelift}: it and the identity are lifts of $p$). It maps lifts of paths to lifts of the same paths, so it
commutes with the action on the fibre: $\tau(\tilde x\cdot g)=\tau(\tilde x)\cdot g$.

\begin{theorem}\label{thm:deck}
Let $p\colon\tilde X\to X$ be a covering with $\tilde X$ path-connected, $G=\pi_1(X,x_0)$, and $H=p_*\pi_1(\tilde X,\tilde x_0)$.
\begin{enumerate}
\item[(a)] If $\tilde x_1=\tilde x_0\cdot g$, then $p_*\pi_1(\tilde X,\tilde x_1)=g^{-1}Hg$.
\item[(b)] There is a deck transformation $\tau$ with $\tau(\tilde x_0)=\tilde x_0\cdot g$ if and only if $g$ lies in the normalizer $N(H)$ of $H$.
\item[(c)] $\operatorname{Deck}(p)\cong N(H)/H$. In particular $p$ is normal if and only if $H$ is a normal subgroup, and then $\operatorname{Deck}(p)\cong G/H$; for a universal covering,
$\operatorname{Deck}(p)\cong G$.
\end{enumerate}
\end{theorem}

\begin{proof}
(a) Let $g=[\gamma]$ and $\tilde\gamma$ the lift of $\gamma$ from $\tilde x_0$, ending at $\tilde x_1$. The isomorphism $\beta_{\tilde\gamma}$ carries $\pi_1(\tilde X,\tilde x_0)$ onto $\pi_1(\tilde X,\tilde x_1)$, and
$p_*\beta_{\tilde\gamma}[\tilde\lambda]=[\bar\gamma\cdot(p\circ\tilde\lambda)\cdot\gamma]=g^{-1}p_*[\tilde\lambda]g$.

(b) By the lifting criterion applied to $f=p$, a lift $\tau$ of $p$ with $\tau(\tilde x_0)=\tilde x_1$ exists if and only if $H\subset p_*\pi_1(\tilde X,\tilde x_1)$, and a lift $\tau'$ with $\tau'(\tilde x_1)=\tilde x_0$ exists
if and only if the reverse inclusion holds. When both exist, $\tau'\circ\tau$ and $\tau\circ\tau'$ are lifts of $p$ fixing a point, hence identities, and $\tau$ is a deck transformation. So $\tau$ exists if and
only if $H=g^{-1}Hg$, that is, $g\in N(H)$. (Inclusion one way suffices for a lift, but a deck transformation needs both.)

(c) For $g\in N(H)$ let $\tau_g$ be the deck transformation with $\tau_g(\tilde x_0)=\tilde x_0\cdot g$. Then $\tau_g\tau_h(\tilde x_0)=\tau_g(\tilde x_0\cdot h)=\tilde x_0\cdot gh$, so $\tau_g\tau_h=\tau_{gh}$ and $g\mapsto\tau_g$ is a
homomorphism. Every deck transformation $\tau$ maps $\tilde x_0$ into the fibre, which is $\tilde x_0\cdot G$ by Proposition~\ref{prop:fibre}, so $\tau=\tau_g$ for some $g$, in $N(H)$ by (b). The kernel is the
stabilizer of $\tilde x_0$, which is $H$. If $H$ is normal, $N(H)=G$ and $\operatorname{Deck}(p)$ reaches every point of the fibre; conversely, transitivity on the fibre means that every $g$
lies in $N(H)$. A universal covering has $H$ trivial.
\end{proof}

\begin{example}\label{ex:computations}
(a) $\R^n\to T^n=\R^n/\Z^n$ is a universal covering whose deck transformations are the translations by $\Z^n$, so $\pi_1(T^n)\cong\Z^n$; for $n=1$ this is Theorem~\ref{thm:pi1circle} again.
(b) For $n\ge2$, $S^n$ is simply connected (Corollary~\ref{cor:spheres} below), so $S^n\to\R P^n$ is universal with deck group $\{\pm\mathrm{id}\}$, and $\pi_1(\R P^n)\cong\Z/2$. (c) For the
$n$-sheeted covering $z\mapsto z^n$ of the circle, $H=n\Z$, which is normal, and the deck group $\Z/n$ consists of the rotations by $n$-th roots of unity.
\end{example}

The analogy with Galois theory (Volume~2, Chapter~10) is exact: the universal covering plays the role of the algebraic closure, the fundamental group the role of the absolute Galois
group, intermediate coverings the role of intermediate fields, and normal coverings the role of normal extensions. It remains to show that every subgroup actually comes from a
covering. A space is \emph{semilocally simply connected} if each point $x$ has a neighbourhood $U$ such that every loop in $U$ at $x$ is nullhomotopic in $X$. Manifolds are (take $U$ a ball in a
chart), and so is every graph (Section~\ref{sec:graphs}).

\begin{theorem}[classification of coverings]\label{thm:galois}
Let $X$ be path-connected, locally path-connected and semilocally simply connected, and $x_0\in X$. For every subgroup $H$ of $\pi_1(X,x_0)$ there is a covering
$p_H\colon(X_H,\tilde x_0)\to(X,x_0)$ with $X_H$ path-connected and $p_{H*}\pi_1(X_H,\tilde x_0)=H$. Two such based coverings with the same subgroup are isomorphic by a unique isomorphism
preserving base points. For $H=\{1\}$ we obtain a universal covering.
\end{theorem}

\begin{proof}
\emph{Construction.} Call an open set \emph{good} if it is path-connected and every loop in it is nullhomotopic in $X$; good sets form a basis of the topology, by the two
hypotheses (a path-connected open subset of a good set is good). Let $X_H$ be the set of classes of paths $\gamma$ starting at $x_0$, for the relation $\gamma\sim\gamma'$ if
$\gamma(1)=\gamma'(1)$ and $[\gamma\cdot\bar\gamma']\in H$ (an equivalence relation because $H$ is a subgroup), and let $p_H\langle\gamma\rangle=\gamma(1)$, where $\langle\gamma\rangle$ denotes the class.
For a good $U$ and $\gamma$ ending in $U$, let $U_{\langle\gamma\rangle}$ be the set of $\langle\gamma\cdot\eta\rangle$ with $\eta$ a path in $U$ from $\gamma(1)$.

(i) $p_H$ maps $U_{\langle\gamma\rangle}$ bijectively onto $U$: onto because $U$ is path-connected; one-to-one because if $\eta,\eta'$ in $U$ end at the same point, $\eta\cdot\bar\eta'$ is nullhomotopic in $X$, so
$\gamma\eta\bar\eta'\bar\gamma\simeq c$ and $\gamma\eta\sim\gamma\eta'$. (ii) If $\langle\gamma'\rangle\in U_{\langle\gamma\rangle}$ then $U_{\langle\gamma'\rangle}=U_{\langle\gamma\rangle}$: if $\gamma'\sim\gamma\eta$, then $\gamma'\mu\sim\gamma(\eta\mu)$, and $\gamma\mu=\gamma\eta\bar\eta\mu\sim\gamma'(\bar\eta\mu)$.
(iii) By (ii), the sets $U_{\langle\gamma\rangle}$ form a basis of a topology on $X_H$: a point $\langle\gamma''\rangle$ of $U_{\langle\gamma\rangle}\cap V_{\langle\gamma'\rangle}$ lies in $W_{\langle\gamma''\rangle}$, for a good $W\subset U\cap V$
containing $\gamma''(1)$, and $W_{\langle\gamma''\rangle}\subset U_{\langle\gamma''\rangle}\cap V_{\langle\gamma''\rangle}=U_{\langle\gamma\rangle}\cap V_{\langle\gamma'\rangle}$. (iv) $p_H$ restricted to $U_{\langle\gamma\rangle}$ is a homeomorphism onto $U$: by (i) and (ii) it carries the basic sets
$V_{\langle\gamma'\rangle}\subset U_{\langle\gamma\rangle}$ (with $V\subset U$ good) onto the good sets $V\subset U$, and conversely. (v) $p_H^{-1}(U)$ is the union of the sets $U_{\langle\gamma\rangle}$, which by (ii) are equal or disjoint:
$U$ is evenly covered, and $p_H$ is a covering.

\emph{Path-connectedness and the subgroup.} For a path $\gamma$ from $x_0$, let $\gamma_t(u)=\gamma(tu)$. The path $t\mapsto\langle\gamma_t\rangle$ in $X_H$ is continuous: near $t_0$, $\gamma$ maps a small interval
into a good $U$, and there $\langle\gamma_t\rangle$ equals $(p_H|_{U_{\langle\gamma_{t_0}\rangle}})^{-1}(\gamma(t))$ (since $\gamma_t$ is $\gamma_{t_0}$ followed by a piece of $\gamma$ in $U$, up to reparametrization). It is therefore
the lift of $\gamma$ from $\tilde x_0=\langle c_{x_0}\rangle$, and it ends at $\langle\gamma\rangle$. So $X_H$ is path-connected, and a loop $\gamma$ at $x_0$ lifts to a loop exactly when $\gamma\sim c_{x_0}$, that is,
$[\gamma]\in H$. By Proposition~\ref{prop:fibre}, $p_{H*}\pi_1(X_H,\tilde x_0)=H$.

\emph{Uniqueness.} If $p_1,p_2$ are based coverings with the same subgroup, the lifting criterion gives lifts $\phi$ of $p_1$ through $p_2$ and $\psi$ of $p_2$ through $p_1$, preserving base points;
$\psi\circ\phi$ and $\phi\circ\psi$ are lifts of $p_1$, $p_2$ fixing a point, hence identities.
\end{proof}

Changing the base point in the fibre conjugates the subgroup (Theorem~\ref{thm:deck}(a)), so connected coverings up to isomorphism, without base points, correspond to conjugacy classes of
subgroups of $\pi_1(X)$. Large subgroups give coverings with few sheets; the whole group gives $X$ itself; the trivial group gives the universal covering, which covers every other
connected covering.

\section{Free products and the theorem of Seifert and van Kampen}

To compute fundamental groups we glue spaces from pieces. The algebra needed is the free product. Let $(G_\alpha)$ be a family of groups. A \emph{reduced word} is a finite
sequence $g_1g_2\cdots g_m$ ($m\ge0$) with each $g_i$ a non-identity element of some $G_{\alpha_i}$ and $\alpha_i\ne\alpha_{i+1}$; let $W$ be the set of reduced words.

\begin{proposition}\label{prop:freeproduct}
There is a group $\mathop{\ast}_\alpha G_\alpha$, the \emph{free product}, whose elements are the reduced words, in which each $G_\alpha$ sits as the words of length at most $1$, and in which the
product of $g_1\cdots g_m$ and $h_1\cdots h_n$ is obtained from the concatenation by repeatedly multiplying adjacent letters from the same group and deleting identities. It has the
universal property: homomorphisms $\varphi_\alpha\colon G_\alpha\to K$ extend uniquely to a homomorphism $\mathop{\ast}_\alpha G_\alpha\to K$.
\end{proposition}

\begin{proof}
The difficulty is associativity, which van der Waerden settled by a trick: make the words act. For $g\in G_\alpha$ define $L_g\colon W\to W$ by: $L_1=\mathrm{id}$; for $g\ne1$ and $w=g_1\cdots g_m$,
$L_gw=gg_1\cdots g_m$ if $g_1\notin G_\alpha$ (or $m=0$), $L_gw=(gg_1)g_2\cdots g_m$ if $g_1\in G_\alpha$ and $gg_1\ne1$, and $L_gw=g_2\cdots g_m$ if $gg_1=1$. A check of the cases shows $L_gL_h=L_{gh}$ for
$g,h$ in the same $G_\alpha$; in particular $L_g$ is a bijection with inverse $L_{g^{-1}}$. Let $P$ be the group of permutations of $W$ generated by all the $L_g$. Every element of $P$ is a product
$L_{g_1}\cdots L_{g_m}$, and using $L_gL_h=L_{gh}$ and $L_1=\mathrm{id}$ we may take $g_1\cdots g_m$ reduced. Applied to the empty word, $L_{g_1}\cdots L_{g_m}$ gives back $g_1\cdots g_m$, so the
reduced word is determined by the permutation: $w\mapsto L_w=L_{g_1}\cdots L_{g_m}$ is a bijection $W\to P$. Transport the group structure of $P$ to $W$; the product in $P$ is the concatenation
followed by the reductions described, and associativity is inherited from composition of maps.

Universal property: on reduced words, define
\[\Phi(g_1\cdots g_m)=\varphi(g_1)\cdots\varphi(g_m).\] The three cases in the definition of $L_g$ show that $\Phi(L_gw)=\varphi(g)\Phi(w)$ for every $g$ and $w$; by
induction on the length of $v$, $\Phi(vw)=\Phi(v)\Phi(w)$. Uniqueness is clear since the $G_\alpha$ generate.
\end{proof}

The \emph{free group} $F(S)$ on a set $S$ is the free product of copies of $\Z$, one for each $s\in S$; its elements are the reduced words in the letters $s^{\pm1}$, and a map from $S$ to a
group $K$ extends uniquely to a homomorphism $F(S)\to K$. We write $F_n$ when $|S|=n$; $F_2$ appeared in Volume~7, Chapter~8. Given homomorphisms $i_1\colon A\to G_1$ and $i_2\colon A\to G_2$, the
\emph{amalgamated product} $G_1\ast_AG_2$ is the quotient of $G_1\ast G_2$ by the normal subgroup $N$ generated by the elements $i_1(a)i_2(a)^{-1}$: it is the largest quotient in which the
two images of $A$ are identified.

\begin{theorem}[Seifert--van Kampen]\label{thm:vankampen}
Let $X=U\cup V$ with $U,V$ open and $U$, $V$, $U\cap V$ path-connected, and $x_0\in U\cap V$. Let $i_U,i_V$ be the maps induced by the inclusions of $U\cap V$, and $j_U,j_V$ those induced by the
inclusions into $X$. The homomorphism $\Phi\colon\pi_1(U,x_0)\ast\pi_1(V,x_0)\to\pi_1(X,x_0)$ extending $j_U$ and $j_V$ is onto, and its kernel is the normal subgroup $N$ generated by the
$i_U(\omega)i_V(\omega)^{-1}$, $\omega\in\pi_1(U\cap V,x_0)$. So
\[\pi_1(X)\cong\pi_1(U)\ast_{\pi_1(U\cap V)}\pi_1(V).\]
\end{theorem}

\begin{proof}
A fact used throughout: every intersection of sets taken from $\{U,V\}$ is $U$, $V$ or $U\cap V$, so it is path-connected and contains $x_0$.

\emph{$\Phi$ is onto.} Let $f$ be a loop at $x_0$. By the Lebesgue number of the cover $\{f^{-1}(U),f^{-1}(V)\}$ of $I$, there is a subdivision $0=s_0<\dots<s_m=1$ with $f([s_{i-1},s_i])\subset A_i\in\{U,V\}$.
Choose a path $g_i$ from $x_0$ to $f(s_i)$ inside $A_i\cap A_{i+1}$ ($g_0$, $g_m$ constant), and let $f_i$ be $f$ restricted to $[s_{i-1},s_i]$, reparametrized. Then
$f\simeq(f_1\bar g_1)(g_1f_2\bar g_2)\cdots(g_{m-1}f_m)$, a product of loops at $x_0$ lying in $U$ or $V$.

\emph{$N$ lies in the kernel}, since $j_Ui_U=j_Vi_V$ (both are induced by the inclusion of $U\cap V$ in $X$).

\emph{The kernel lies in $N$.} Call a \emph{factorization} of a class $[f]\in\pi_1(X,x_0)$ a formal product $[f_1]\cdots[f_k]$ in which each $f_i$ is a loop at $x_0$ in some $A_i\in\{U,V\}$, $[f_i]$ is its
class in $\pi_1(A_i,x_0)$, and $f\simeq f_1\cdots f_k$ in $X$. A factorization is a word in the free product, mapped to $[f]$ by $\Phi$. Consider two moves: (1) combining two adjacent terms in the
same $\pi_1(A)$ into their product, or the reverse; (2) when $f_i$ lies in $U\cap V$, regarding $[f_i]$ in $\pi_1(V)$ instead of $\pi_1(U)$, or the reverse. Move (1) does not change the element of
the free product, and move (2) does not change its image in $Q=(\pi_1(U)\ast\pi_1(V))/N$. We show that any two factorizations of the same $[f]$ are related by these moves. Applied to a
word in the kernel and to the empty word, both factorizations of $[c_{x_0}]$, this shows that the word lies in $N$.

Let $f_1\cdots f_k$ and $f'_1\cdots f'_l$ be factorizations of $[f]$ and $F\colon I\times I\to X$ a homotopy from $f_1\cdots f_k$ to $f'_1\cdots f'_l$ fixing the end points. Subdivide the square into rectangles $R$, by
subdivisions of the bottom edge refining the points $\frac ik$ and of the top edge refining the points $\frac jl$ (use a common vertical subdivision containing all of these), so fine that each
$F(R)$ lies in a set $A(R)\in\{U,V\}$. For each vertex $v$ of the grid choose a path $g_v$ from $x_0$ to $F(v)$ lying in every $A(R)$ for the rectangles $R$ containing $v$, and moreover in
$A_i$ if $v$ is on the bottom edge inside $[\frac{i-1}k,\frac ik]$, or in $A'_j$ if $v$ is on the top edge inside $[\frac{j-1}l,\frac jl]$; this is possible by the fact above, since $F(v)$ lies in each of these sets.
At vertices with $F(v)=x_0$ (the four sides of the square, and the points $\frac ik$, $\frac jl$) take $g_v$ constant. For an edge $\sigma$ of the grid from $v$ to $w$, the loop
$L(\sigma)=g_v\cdot(F\circ\sigma)\cdot\bar g_w$ lies in $A(R)$ for every rectangle $R$ containing $\sigma$, so it defines an element of $\pi_1(A(R))$; if two rectangles containing $\sigma$ have different sets, $L(\sigma)$ lies
in $U\cap V$, and the two choices differ by move (2).

For an edge path $\gamma$ in the grid from $(0,0)$ to $(1,1)$, made of edges $\sigma_1,\dots,\sigma_q$, the product $L(\sigma_1)\cdots L(\sigma_q)$ is homotopic to $F\circ\gamma$ (the paths $g_v$ cancel in pairs),
so it is a factorization of $[f]$ (the homotopy $F$ carries $F\circ\gamma$ to the bottom-right path, which is $f_1\cdots f_k$ followed by the constant right side, by Lemma~\ref{lem:convexhomotopy}). Let $\gamma_0$
run along the bottom edge and up the right side, and push it across the rectangles one at a time, row by row from the bottom left, until it runs up the left side and along the top.
Crossing a rectangle $R$ replaces its bottom and right edges by its left and top edges. All four loops lie in $A(R)$, and in $A(R)$ the product along the bottom and right is homotopic to the
product along the left and top, by the homotopy $F$ restricted to $R$; so this step is a move (1), after moves (2) to view the four terms in $\pi_1(A(R))$. Finally, the factorization of $\gamma_0$
is equivalent to $[f_1]\cdots[f_k]$: the right-side terms are constant loops (the identity), and the bottom edges inside $[\frac{i-1}k,\frac ik]$ give loops in $A_i$, whose product, viewed in $\pi_1(A_i)$
by moves (2) and combined by moves (1), is $[f_i]$. Likewise the last path gives $[f'_1]\cdots[f'_l]$.
\end{proof}

\begin{corollary}\label{cor:spheres}
For $n\ge2$ the sphere $S^n$ is simply connected.
\end{corollary}

\begin{proof}
Let $U=S^n\setminus\{N\}$ and $V=S^n\setminus\{S\}$, the complements of the poles. Stereographic projection (Exercise~\ref{ex:stereographic}) maps each homeomorphically onto $\R^n$, which is simply connected,
and maps $U\cap V$ onto $\R^n\setminus\{0\}$, which is path-connected for $n\ge2$. By Theorem~\ref{thm:vankampen}, $\pi_1(S^n)$ is a quotient of the free product of two trivial groups.
\end{proof}

For $n=1$ the hypothesis fails, $U\cap V$ being two arcs, and indeed $\pi_1(S^1)=\Z$. Attaching a disc is the other basic gluing. Let $X$ be a space and $\varphi\colon S^1\to X$ a map; $X\cup_\varphi D^2$
is the quotient of the disjoint union of $X$ and the closed disc $D^2$ in which each $z\in S^1$ is identified with $\varphi(z)$.

\begin{proposition}\label{prop:attachdisc}
If $X$ is path-connected, $\pi_1(X\cup_\varphi D^2)\cong\pi_1(X)/\langle\langle[\varphi]\rangle\rangle$, the quotient by the normal subgroup generated by the class of the loop $\varphi$ (moved to the base point
by any path).
\end{proposition}

\begin{proof}
Let $Y=X\cup_\varphi D^2$, $U=Y\setminus\{0\}$ (the centre of the disc removed) and $V$ the open disc. Pushing radially outwards, $U$ deformation retracts onto $X$; $V$ is simply connected;
$U\cap V$ is a punctured disc, whose fundamental group is $\Z$, generated by a circle of radius $\frac12$. Take the base point $y_0=\frac12$ in $U\cap V$. Under the retraction, that circle becomes,
up to the change of base point along the radius from $\frac12$ to $1$, the loop $\varphi$. Theorem~\ref{thm:vankampen} gives $\pi_1(Y,y_0)\cong\pi_1(U,y_0)/N$, where $N$ is the normal subgroup generated by
the image of $\pi_1(U\cap V)$, that is, by the class of $\varphi$ transported to $y_0$.
\end{proof}

\begin{example}\label{ex:surfaces}
Glue the sides of a $4g$-gon in the pattern
\[a_1b_1a_1^{-1}b_1^{-1}\cdots a_gb_ga_g^{-1}b_g^{-1}:\]
the result is the closed orientable surface $\Sigma_g$ of genus $g$ (for $g=1$, the square with opposite
sides glued, the torus). All vertices of the polygon become one point, and the edges become $2g$ circles through it, a wedge of $2g$ circles; the inside of the polygon is a disc attached along
the word. By Proposition~\ref{prop:attachdisc} and Theorem~\ref{thm:wedge} below,
\[\pi_1(\Sigma_g)\cong\bigl\langle a_1,b_1,\dots,a_g,b_g\mid[a_1,b_1]\cdots[a_g,b_g]\bigr\rangle,\qquad[a,b]=aba^{-1}b^{-1},\]
the quotient of the free group by the normal subgroup generated by the relator. Making the group abelian kills the relator and gives $\Z^{2g}$, so surfaces of different genus have
non-isomorphic fundamental groups and are not homeomorphic, nor even homotopy equivalent. For $g=1$ we recover $\Z^2$. The projective plane is a disc with antipodal boundary points identified,
a circle $a$ with a disc attached along $a^2$: $\pi_1(\R P^2)=\Z/2$, as in Example~\ref{ex:computations}.
\end{example}

\section{Graphs and free groups}\label{sec:graphs}

A \emph{graph} is a space obtained from a set of vertices and a set of edges, each edge a copy of $[0,1]$ whose ends are attached to vertices (loops and multiple edges allowed). It is
\emph{locally finite} if each vertex meets finitely many edge ends. We use the path metric in which each edge has length $1$; a locally finite graph is then locally path-connected, and it is
semilocally simply connected since the open star of a vertex (the vertex with the open half-edges at it) contains no loop that is not nullhomotopic in the star (it deformation retracts onto
the vertex). A \emph{tree} is a connected graph without embedded circles; a \emph{spanning tree} of a connected graph is a subgraph that is a tree and contains all the vertices.

\begin{lemma}\label{lem:trees}
\begin{enumerate}
\item[(a)] If $v$ is a vertex of a graph $Y$ lying on a single edge end, of an edge $e$ that is not a loop, then $Y$ deformation retracts onto $Y\setminus(\{v\}\cup e^\circ)$, where $e^\circ$ is the open edge.
\item[(b)] A finite tree with $V$ vertices has $V-1$ edges and deformation retracts onto any of its vertices; a finite tree with one extra edge is homotopy equivalent to a circle.
\item[(c)] Every finite connected graph has a spanning tree.
\end{enumerate}
\end{lemma}

\begin{proof}
(a) Slide $e$ into its other end point along itself. (b) A finite tree with at least two vertices has a vertex on a single edge end: follow a path without backtracking as long as possible;
since there is no circle, it never revisits a vertex, so it stops, at such a vertex. Remove it by (a) and induct. With one extra edge, remove such end vertices as long as there are any; what
remains is a connected graph in which every vertex has two edge ends, a circle (if no vertex lies on a single edge end and the graph has as many edges as vertices, every vertex lies on exactly
two edge ends). (c) Remove edges that lie on embedded circles, one at a time, while there are any; removing such an edge keeps the graph connected.
\end{proof}

\begin{theorem}\label{thm:wedge}
Let $X$ be a finite connected graph with $V$ vertices and $E$ edges, $T$ a spanning tree, and $x_0$ a vertex. For each edge $e$ not in $T$, oriented from $u$ to $w$, let $\gamma_e$ be the loop that goes in
$T$ from $x_0$ to $u$, along $e$, and back in $T$ to $x_0$. Then $\pi_1(X,x_0)$ is the free group on the classes $[\gamma_e]$, of rank $E-V+1$. In particular $\pi_1$ of a wedge of $n$ circles (one
vertex, $n$ loop edges) is $F_n$.
\end{theorem}

\begin{proof}
Induction on the number of edges outside $T$. If there are none, $X=T$ is simply connected by Lemma~\ref{lem:trees}(b). Otherwise pick $e\notin T$ with midpoint $m$, and let $X'=X\setminus e^\circ$, a
connected graph with spanning tree $T$. Let $U=X\setminus\{m\}$ and $V$ the union of $e^\circ$ and the open $\frac13$-neighbourhood of $T$. Then $U$ deformation retracts onto $X'$ (push the two
halves of $e$ to its ends); $V$ deformation retracts onto $T\cup e$, a circle up to homotopy, with $\pi_1$ generated by $[\gamma_e]$ (under the retraction of Lemma~\ref{lem:trees}(b), $\gamma_e$ goes once around the circle); and $U\cap V$ deformation retracts onto $T$ with two half-edges
attached, which retracts further onto $T$, so it is simply connected. All three are path-connected. Theorem~\ref{thm:vankampen} gives $\pi_1(X)\cong\pi_1(X')\ast\Z$, with the new factor generated
by $[\gamma_e]$, and induction finishes. The rank is the number of edges outside $T$, which is $E-(V-1)$.
\end{proof}

\begin{figure}[ht]
\centering
\fig{ch05-coverings}
\caption{The figure eight $R_2$, with loops $a$ (orange) and $b$ (blue), and its three connected two-sheeted coverings. Each vertex has exactly one outgoing and one incoming edge of each
colour; reading edge labels along loops at the lower vertex gives the subgroup of index $2$: $\langle a,\,b^2,\,bab^{-1}\rangle$, $\langle b,\,a^2,\,aba^{-1}\rangle$ and $\langle a^2,\,ab,\,b^2\rangle$, each free of rank $3$.}
\label{fig:coverings}
\end{figure}

Coverings of graphs are graphs, which can be drawn. For a subgroup $H$ of $F_n=\pi_1(R_n)$, where $R_n$ is the wedge of $n$ circles $a_1,\dots,a_n$, the \emph{Schreier coset graph} $X_H$ has the right
cosets $Hg$ as vertices and, for each coset $Hg$ and each $i$, an edge labelled $a_i$ from $Hg$ to $Hga_i$. Mapping each vertex to the vertex of $R_n$ and each edge labelled $a_i$ onto the circle
$a_i$ is a covering: at each vertex there is exactly one outgoing and one incoming edge with each label, because $g\mapsto ga_i$ permutes the cosets, so a small star around the vertex of $R_n$ is
evenly covered. The lift from the vertex $H$ of the loop spelled by a word $w$ ends at $Hw$; so by Proposition~\ref{prop:fibre} the subgroup of this covering is $\{w:Hw=H\}=H$. This is
Theorem~\ref{thm:galois} made concrete (Figure~\ref{fig:coverings}).

\begin{theorem}[Nielsen--Schreier]\label{thm:nielsenschreier}
Every subgroup $H$ of the free group $F_n$ is free. If $H$ has finite index $k$, its rank is
\[1+k(n-1).\]
\end{theorem}

\begin{proof}
$X_H$ is a connected covering of $R_n$, and $p_*$ is injective (Corollary~\ref{cor:liftendpoints}), so $H\cong\pi_1(X_H)$. If $k$ is finite, $X_H$ is a finite graph with $k$ vertices and $nk$ edges, and
Theorem~\ref{thm:wedge} gives a free group of rank $nk-k+1$.

If $k$ is infinite, $X_H$ is a countable, locally finite, connected graph (the cosets are countable since $F_n$ is). Let $v_0=H$ and, for each other vertex $v$, choose one edge joining $v$ to a vertex
nearer to $v_0$ (in the path metric). The chosen edges form a spanning tree $T$: following them from any vertex reaches $v_0$, and on an embedded circle made of chosen edges, the vertex
farthest from $v_0$ would have chosen both of its circle edges, since each chosen edge joins vertices whose distances to $v_0$ differ by $1$ and was chosen by the farther one. Define $\gamma_e$ for $e\notin T$ as before. Every compact set $K\subset X_H$ lies in a finite subgraph, since the open balls of radius $1$ around the vertices cover $K$, finitely many of them suffice, and each meets finitely many edges.
Given finitely many edges and a compact set, the union $Y$ of the edges, the finite subgraph containing the compact set, and the paths in $T$ from $v_0$ to all their vertices is a finite
connected graph in which $T\cap Y$ is a spanning tree. A loop at $v_0$ lies in such a $Y$, where it is a product of the $[\gamma_e]$ with $e$ in $Y\setminus T$, by Theorem~\ref{thm:wedge}; so these classes
generate. If a nonempty reduced word in the $[\gamma_e]$ were trivial in $\pi_1(X_H)$, the image of a nullhomotopy would lie in such a $Y$ containing the edges involved, contradicting the freeness
of $\pi_1(Y)$ on its $\gamma_e$. So $\pi_1(X_H)$ is free on the $[\gamma_e]$, $e\notin T$.
\end{proof}

\begin{beauty}
Nielsen proved in 1921 that finitely generated subgroups of free groups are free, by a long cancellation argument on words; Schreier proved the general case and the index formula in 1927,
with coset representatives. With coverings the theorem becomes almost obvious: the subgroup is the fundamental group of a covering, the covering of a graph is a graph, and the
fundamental group of a graph is free. The index formula is an Euler characteristic count, $\chi(X_H)=k\,\chi(R_n)$, that is, $1-\operatorname{rank}H=k(1-n)$. A subgroup of index $2$ in $F_2$ already needs
$3$ generators, and a subgroup of index $k$ needs $k+1$: subgroups of free groups are bigger than the group, a phenomenon unknown for vector spaces.
\end{beauty}

\begin{remark}
The theorem holds for free groups of any rank, by the same proof with Zorn's lemma for spanning trees. The fundamental group of a surface is not free, but Example~\ref{ex:surfaces} with
Theorem~\ref{thm:galois} still gives a dictionary between its subgroups and the coverings of the surface; a covering of degree $k$ of $\Sigma_g$ is a surface $\Sigma_h$ with $2-2h=k(2-2g)$, the same Euler
characteristic count, which Chapter~6 justifies.
\end{remark}

\section{Exercises}

\begin{exercise}[$\star$]\label{ex:fta}
\emph{The fundamental theorem of algebra.} Let $P(z)=z^n+a_{n-1}z^{n-1}+\dots+a_0$ with $n\ge1$, and suppose $P$ has no root. (a) Show that $\gamma_r(t)=\frac{P(re^{2\pi it})/P(r)}{|P(re^{2\pi it})/P(r)|}$ is a loop at
$1$ in $S^1$ whose class does not depend on $r\ge0$. (b) For $r=0$ it is constant; for $r$ large, compare it with $\omega_n$ and conclude.
\end{exercise}

\begin{exercise}[$\star\star$]\label{ex:borsukulam2}
\emph{Borsuk--Ulam in dimension $2$.} (a) Show that there is no map $g\colon S^2\to S^1$ with $g(-x)=-g(x)$ for all $x$. (Restrict $g$ to the equator, lift to $\R$, and show the loop has odd
winding number.) (b) Deduce that every map $f\colon S^2\to\R^2$ takes the same value at two antipodal points: at every moment there are two antipodal points on the Earth with the same
temperature and the same pressure. (c) Deduce that $S^2$ is not homeomorphic to a subset of $\R^2$.
\end{exercise}

\begin{exercise}[$\star$]\label{ex:figureeight}
\emph{Coverings of the figure eight.} (a) Show that the subgroups of index $2$ of $F_2=\langle a,b\rangle$ are the kernels of the three nonzero homomorphisms $F_2\to\Z/2$, and match them with
Figure~\ref{fig:coverings}. (b) Find free bases of each, check the rank $3$ predicted by Theorem~\ref{thm:nielsenschreier}, and decide which of these coverings are normal.
\end{exercise}

\begin{exercise}[$\star\star$]\label{ex:commutator}
\emph{The commutator subgroup.} Let $K$ be the kernel of the homomorphism $F_2\to\Z^2$ sending $a,b$ to the basis vectors. (a) Show that the corresponding covering of the figure eight is the
grid graph $\Z^2\subset\R^2$ (vertices at integer points, horizontal edges labelled $a$, vertical edges labelled $b$). (b) Deduce that $K$ is free of infinite rank, and describe a free basis. (c) Show that $K$ is the subgroup generated by all commutators $[g,h]$.
\end{exercise}

\begin{exercise}[$\star\star$]\label{ex:picture}
\emph{Hanging a picture.} A picture hangs on a string looped around two nails in a wall. Model the wall without the nails as $X=\R^2\setminus\{p,q\}$ and the string as a loop at a point below the
nails. (a) Show that $X$ deformation retracts onto a figure eight, so $\pi_1(X)=F_2$, freely generated by loops $a$ and $b$ around $p$ and $q$. (b) The picture falls when the string becomes
nullhomotopic. Pulling out the nail $q$ maps $\pi_1(X)\to\pi_1(\R^2\setminus\{p\})=\Z$ by $a\mapsto1$, $b\mapsto0$. Show that the string $aba^{-1}b^{-1}$ keeps the picture up with both nails, but that it falls
if either nail is removed. (c) Find a string for three nails that falls if any one of them is removed.
\end{exercise}

\begin{exercise}[$\star$]\label{ex:rp3}
\emph{Projective space and the sphere.} (a) Using Example~\ref{ex:computations}, show that every loop in $\R P^3$ is either nullhomotopic or homotopic to the image of a half great circle of $S^3$,
and that this one is not nullhomotopic though its square is. (Chapter~7 identifies $\R P^3$ with the rotation group $SO(3)$; the loop becomes a full turn about an axis, and its square the double turn
of the plate trick of Volume~2.) (b) Show that every map $S^2\to S^1$, and every map $S^2\to T^2$, is homotopic to a constant map.
\end{exercise}

\section{Hints}

\begin{hint}{ex:fta}
For (a), a homotopy is given by letting $r$ vary. For (b), divide by $r^n$: for $r$ large, $P(re^{2\pi it})$ stays close to $r^ne^{2\pi int}$, and the dog-on-a-leash principle applies.
\end{hint}

\begin{hint}{ex:borsukulam2}
For (a): if $\tilde h$ lifts $h(t)=g(\cos\pi t,\sin\pi t,0)$, oddness gives $\tilde h(t+1)=\tilde h(t)+\text{half an odd integer}$. For (b), apply (a) to $\frac{f(x)-f(-x)}{|f(x)-f(-x)|}$.
\end{hint}

\begin{hint}{ex:figureeight}
A subgroup of index $2$ is normal, and the quotient is $\Z/2$. Use the spanning tree made of one edge between the two vertices.
\end{hint}

\begin{hint}{ex:commutator}
For (a), the vertices of the covering are the cosets $Kg$, which correspond to the images of $g$ in $\Z^2$. For (b), choose the spanning tree formed by the horizontal axis and all vertical lines;
then remove the edges outside the tree one square at a time. For (c), the quotient by the commutators is abelian.
\end{hint}

\begin{hint}{ex:picture}
For (c), a commutator of a commutator and the third nail: $[[a,b],c]$.
\end{hint}

\begin{hint}{ex:rp3}
For (a), the lift of a loop to $S^3$ ends at the starting point or at its antipode. For (b), use the lifting criterion with the universal coverings $\R\to S^1$ and $\R^2\to T^2$.
\end{hint}

\section{Solutions}

\begin{solution}{ex:fta}
\idea A loop that cannot be undone is not nullhomotopic: if $P$ had no root, the image of a huge circle would wind $n$ times around $0$ and still be shrinkable through the images of smaller circles.
\writeup (a) $H(t,s)=\gamma_{s}(t)$ for $s$ ranging over $[0,r]$ is continuous because $P$ does not vanish, and each $\gamma_s$ is a loop at $1$ in $S^1$. So all $\gamma_r$ are homotopic to $\gamma_0$, which is constant.
(b) For $|z|=r$ with $r>1+|a_{n-1}|+\dots+|a_0|$, $|P(z)-z^n|\le(|a_{n-1}|+\dots+|a_0|)r^{n-1}<r^n=|z^n|$, so by the dog-on-a-leash principle (Volume~7, Proposition~9.3(a)) $t\mapsto P(re^{2\pi it})$ has winding
number $n$ around $0$. Dividing by the positive number $|P(re^{2\pi it})|$ and by the constant $P(r)$ moves the loop within $\C^*$ without changing its winding number (straight-line homotopy to
$1$ of the modulus, and a path from $P(r)$ to $1$ in $\C^*$), so $[\gamma_r]=n\ne0$ in $\pi_1(S^1)=\Z$, contradicting (a). So $P$ has a root.
\end{solution}

\begin{solution}{ex:borsukulam2}
\idea An odd map sends half the equator to a path from a point to its antipode, which turns an odd number of half-turns; doubled, the equator turns an odd number of times, so it cannot
bound the disc that a hemisphere provides.
\writeup (a) Let $h(t)=g(\cos\pi t,\sin\pi t,0)$ on $[0,2]$, a loop (the equator), with $h(t+1)=-h(t)$ by oddness. Lift it, $h=e^{2\pi i\tilde h}$. Then $e^{2\pi i\tilde h(t+1)}=-e^{2\pi i\tilde h(t)}$, so
$\tilde h(t+1)-\tilde h(t)$ is in $\frac12+\Z$; being continuous, it is a constant $q+\frac12$. The winding number of the equator loop is $\tilde h(2)-\tilde h(0)=2q+1$, odd. But the equator is
nullhomotopic in $S^2$ (Corollary~\ref{cor:spheres}), so its image under $g$ is nullhomotopic, of winding number $0$. Contradiction. (b) If $f(x)\ne f(-x)$ for all $x$, then
$g(x)=\frac{f(x)-f(-x)}{|f(x)-f(-x)|}$ is an odd map $S^2\to S^1$. Temperature and pressure form a map $S^2\to\R^2$, assuming they vary continuously. (c) A one-to-one map $S^2\to\R^2$
would contradict (b).
\end{solution}

\begin{solution}{ex:figureeight}
\idea A covering with two sheets is a pair of vertices with the two $a$-edges and the two $b$-edges each either looping or swapping.
\writeup (a) A subgroup of index $2$ contains the square of every element and is normal (its left and right cosets coincide, being $H$ and its complement), so it is the kernel of a homomorphism onto
$\Z/2$. Such a homomorphism is determined by the images of $a$ and $b$, not both zero: $(1,0)$, $(0,1)$, $(1,1)$. In the corresponding coverings, $a$ swaps the vertices exactly when $a\mapsto1$.
(b) Take the tree formed by one edge joining the two vertices. For $\ker(a\mapsto0,b\mapsto1)$: the tree is a $b$-edge, the other edges are the two $a$-loops and the other $b$-edge, giving the basis
$a$, $bab^{-1}$, $b^2$. Symmetrically $b$, $aba^{-1}$, $a^2$ for the second. For the third, the tree is an $a$-edge, and the other edges give $a^2$, $ab$ and $ba^{-1}$; replacing $ba^{-1}$ by $(ba^{-1})(ab)=b^2$ gives the generators $a^2$, $ab$, $b^2$ of the caption.
Rank $3=1+2(2-1)$ in each case. All three are normal, as subgroups of index $2$, and the deck transformation swaps the two vertices.
\end{solution}

\begin{solution}{ex:commutator}
\idea The covering is the graph on which you can walk only knowing how many $a$'s and $b$'s, with signs, you have read: the integer grid.
\writeup (a) The homomorphism $\pi\colon F_2\to\Z^2$ is onto with kernel $K$, so the cosets $Kg$ correspond to the points $\pi(g)\in\Z^2$, and $Kg\cdot a$ corresponds to $\pi(g)+(1,0)$: the coset graph is the grid,
with $a$-edges horizontal and $b$-edges vertical. (b) The index is infinite, so by Theorem~\ref{thm:nielsenschreier} $K$ is free, on the loops $\gamma_e$ for the edges outside a spanning tree. With the
tree made of the $x$-axis and all vertical lines, the edges outside are the horizontal edges off the axis; for such an edge $e$ from $(i,j)$ to $(i+1,j)$, $\gamma_e$ walks along the axis to $(i,0)$, up to $e$, across it, down to $(i+1,0)$
and back along the axis, which is the word $a^ib^ja\,b^{-j}a^{-i-1}$. There are infinitely many, so $K$ has infinite rank. (c) $\Z^2$ is abelian, so every
commutator lies in $K$. Conversely, modulo the subgroup $C$ generated by commutators (which is normal), $F_2/C$ is abelian and generated by the images of $a$ and $b$, so every element of
$F_2/C$ is $a^mb^n$ mod $C$; if $g\in K$, $\pi(g)=(m,n)=0$ and $g\in C$.
\end{solution}

\begin{solution}{ex:picture}
\idea The string is a word in $a^{\pm1}$, $b^{\pm1}$; removing a nail erases its letter; the picture falls exactly when what remains cancels.
\writeup (a) $X$ deformation retracts onto the union of two circles around $p$ and $q$ joined by a segment, which is homotopy equivalent to a figure eight (contract the segment, a tree); so
$\pi_1(X)=F_2$ by Theorem~\ref{thm:wedge}. (b) $aba^{-1}b^{-1}$ is a nonempty reduced word, so it is nontrivial in $F_2$. Removing $q$ sends it to $1+0-1-0=0$ in $\Z$: nullhomotopic. Removing $p$
likewise. (c) Let $w=[a,b]$ and take $[w,c]=wcw^{-1}c^{-1}$ (written out, a word of length $10$, reduced and nonempty, so nontrivial). Removing $c$ gives $ww^{-1}=1$; removing $a$ or $b$ makes $w$
trivial, and then $[w,c]=cc^{-1}=1$. The idea extends to any number of nails by nested commutators; Demaine and his coauthors (2014) found shorter strings.
\end{solution}

\begin{solution}{ex:rp3}
\idea Loops in $\R P^3$ lift to paths in $S^3$ that end where they started or at the antipode; the first kind shrinks, the second kind is one nontrivial class.
\writeup (a) $S^3\to\R P^3$ is a universal covering with deck group $\{\pm1\}$, so $\pi_1(\R P^3)\cong\Z/2$ by Theorem~\ref{thm:deck}. The class of a loop is nontrivial exactly when its lift from $x$ ends at $-x$
(Proposition~\ref{prop:fibre}). The image of a half great circle from $x$ to $-x$ is such a loop; it is not nullhomotopic, and its square is trivial because $\Z/2$ has exponent $2$. Every other nontrivial
loop is homotopic to it, as the group has only one nontrivial element. (b) $S^2$ is simply connected (Corollary~\ref{cor:spheres}), so by Theorem~\ref{thm:liftcrit} a map $f\colon S^2\to S^1$ lifts to
$\tilde f\colon S^2\to\R$. The homotopy $(1-s)\tilde f$ in $\R$, pushed down, joins $f$ to a constant map. The same works with $\R^2\to T^2$.
\end{solution}

\section*{Sources and further reading}
\addcontentsline{toc}{section}{Sources and further reading}

H.~Poincar\'e, ``Analysis situs'', \emph{J.\ \'Ec.\ Polytech.}\ (2) 1 (1895), 1--121, and ``Cinqui\`eme compl\'ement \`a l'Analysis situs'', \emph{Rend.\ Circ.\ Mat.\ Palermo} 18 (1904), 45--110; both are
translated in H.~Poincar\'e, \emph{Papers on Topology: Analysis Situs and Its Five Supplements}, translated by J.~Stillwell (AMS and LMS, 2010), whose introduction tells the story of the corrections. On the
solution: G.~Perelman, arXiv preprints math/0211159, math/0303109 and math/0307245 (2002--2003); J.~Morgan and G.~Tian, \emph{Ricci Flow and the Poincar\'e Conjecture} (AMS and Clay Mathematics
Institute, 2007); the Clay Mathematics Institute's announcement of 18~March~2010; J.~Milnor, ``Towards the Poincar\'e conjecture and the classification of 3-manifolds'', \emph{Notices AMS} 50 (2003), 1226--1233.

H.~Seifert, ``Kon\-struk\-tion drei\-di\-men\-sio\-na\-ler ge\-schlos\-se\-ner R\"aume'', \emph{Ber.\ S\"achs.\ Akad.\ Wiss.} 83 (1931), 26--66; E.~R.~van Kampen, ``On the connection between the fundamental groups of
some related spaces'', \emph{Amer.\ J.\ Math.} 55 (1933), 261--267. J.~Nielsen, ``Om regning med ikke-kommutative faktorer og dens anvendelse i gruppeteorien'', \emph{Mat.\ Tidsskr.\ B} (1921), 77--94;
O.~Schreier, ``Die Untergruppen der freien Gruppen'', \emph{Abh.\ Math.\ Sem.\ Hamburg} 5 (1927), 161--183. E.~D.~Demaine, M.~L.~Demaine, Y.~N.~Minsky, J.~S.~B.~Mitchell, R.~L.~Rivest and M.~P\v atra\c scu,
``Picture-hanging puzzles'', \emph{Theory Comput.\ Syst.} 54 (2014), 531--550.

The proofs follow A.~Hatcher, \emph{Algebraic Topology} (Cambridge University Press, 2002), Sections~1.1--1.3 and~1.A, where the proof of van Kampen's theorem is given for any number of open sets,
and W.~Massey, \emph{Algebraic Topology: An Introduction} (Harcourt, Brace and World, 1967). The proof of Proposition~\ref{prop:freeproduct} is van der Waerden's, from ``Free products of groups'',
\emph{Amer.\ J.\ Math.} 70 (1948), 527--528. For the combinatorial side, J.~Stallings, ``Topology of finite graphs'', \emph{Invent.\ Math.} 71 (1983), 551--565.
